\section{Conclusion}\label{conclusion}

This work evaluated On-Policy Self-Distillation (OPSD) for visual answering using question-grounded
image crops as privileged teacher context. OPSD was compared with conventional Supervised Fine-Tuning (SFT)
across three VLMs (SmolVLM, BLIP and MiniCPM) using the GQA dataset.

The results showed that SFT substantially outperformed OPSD form SmolVLM and BLIP, while the two
methods obtained near-identical BERTScore-F1 performance for MiniCPM. This indicates that OPSD was not 
generally more effective than SFT under the experimental conditions considered, but its relative
effectiveness varied across VLM architectures. The weak crop-teacher advantage observed before training
suggests that the privileged context did not consistently provide a sufficient stronger supervisory signal.

Qualitative analysis further showed that crop-based privileged supervision can be beneficial when relevant
evidence is local and surrounded by distracting visual content, but can be harmful when broader
scene context is required. Overall, the findings suggest that OPSD is highly dependent on the quality
of the privileged context and may be better suited to domains where training-time information provides 
a clear and reliable advantage over the standard inference input.

